\documentclass[11pt]{amsart}
\usepackage{amssymb,amsmath,amsthm}
\usepackage[margin=1.1in]{geometry}
\usepackage[hidelinks]{hyperref}
\usepackage{enumitem}
\usepackage{booktabs}
\setlist{nosep}
\setlist[itemize]{leftmargin=*}

\newtheorem{theorem}{Theorem}
\newtheorem{corollary}[theorem]{Corollary}
\newtheorem{lemma}[theorem]{Lemma}
\newtheorem{proposition}[theorem]{Proposition}
\newtheorem{claim}[theorem]{Claim}
\theoremstyle{definition}
\newtheorem{definition}[theorem]{Definition}
\newtheorem{observation}[theorem]{Observation}
\theoremstyle{remark}
\newtheorem{remark}[theorem]{Remark}

\newcommand{\w}{\omega}
\newcommand{\p}[1]{(#1)}
\newcommand{\cut}[2]{#1\,@\,#2}
\newcommand{\PA}{\mathbf{A}}
\newcommand{\PB}{\mathbf{B}}

\title[Chomp on $3\times3\times\omega$]{Chomp on $3\times 3\times\omega$: a P-position with four
infinite columns,\\ and the small positions around it}
\author{Charlie Reams}
\date{\today}

\begin{document}
\begin{abstract}
Gale asked whether the first player wins Chomp on the box $3\times3\times\omega$. A position
below this box is given by nine column heights in $\mathbb N\cup\{\omega\}$. Positions with one
infinite column are first-player wins, and those with two infinite columns can be decided by a
periodicity argument; beyond that nothing was known. We prove that the position $\PA$ with the
$2\times2$ block of columns infinite, two columns of height $1$ and one of height $2$ is a
P-position (a second-player win). The proof is computer-assisted: a set of finite positions in
fourteen linear families is shown to be closed under replies by an SMT solver, and the one
family of options whose replies are irregular is handled by a strategy-stealing argument.
The theorem is also checked in Lean~4 from the rules of the game.
We then describe exactly how the P-positions with two infinite columns below a position such as
$\PA$ are arranged (a greedy matching), use it with a certified solver to show that a second
candidate $\PB$ is an N-position by a reply at height $31$, and classify all $2177$ positions
with three or four infinite columns and finite columns of height at most $4$: $\PA$ and its
transpose are the only P-positions. All $825$ positions with the cross of five columns infinite
and corner columns of height at most $8$ have a finite winning reply. Finally we report a
regularity in the data: a chain of cuts that has no P member up to height $u$ gets one at
height $u$ with probability close to $1/u$, and a simple model that reproduces it. The box $3\times3\times\omega$ itself remains
undecided.
\end{abstract}
\maketitle

\section{Introduction}

Chomp is played on a partially ordered set with a least element, the poison. A move removes
an element together with everything above it, and the player who takes the poison loses. A
position is a \emph{P-position} if the player to move loses and an \emph{N-position} otherwise.
Gale \cite{Gale} asked who wins on $\mathbb N^3$ and proposed the box $3\times3\times\omega$,
that is $\{0,1,2\}^2\times\mathbb N$, as a possibly simpler question. The question about
$\mathbb N^3$ is listed as open in the current collection of unsolved problems \cite[D1]{Now}. Huddleston and Shurman
\cite{HS} developed the theory of Chomp on boxes with ordinal sides and proved, among other
things, that $2\times\omega$, $3\times\omega^\omega$ and $2\times2\times\omega^3$ are P-positions.

A position below $3\times3\times\omega$ consists of nine columns over the cells of the $3\times3$
base, of heights in $\mathbb N\cup\{\omega\}$. What is known sorts by the number of infinite
columns.
\begin{itemize}
\item \emph{One} (the poison column): every such position is an N-position, by the Fundamental
Theorem of \cite{HS,Gale}.
\item \emph{Two}: Huddleston and Shurman state, without proof, that the set of P-positions
is recursive when the infinite part is two columns wide \cite[p.~211]{HS}. The mechanism is
a periodicity of the kind that Byrnes \cite{By} proved for finite positions with two long
chains. We use a solver that produces a certificate (Section~\ref{sec:exact}).
\item \emph{Three or more}: no general method is known, and to our knowledge no P-position was
known.
\end{itemize}
The box $3\times3\times\omega$ is a P-position if and only if none of its options is one, and
its options have between three and eight infinite columns.

\subsection*{Results}
Positions are written by their column heights $a_{xy}$ in the order
\[
(a_{00},a_{01},a_{02}\mid a_{10},a_{11},a_{12}\mid a_{20},a_{21},a_{22}).
\]

\begin{theorem}\label{thm:A}
$\PA=(\w,\w,1\mid\w,\w,1\mid 2,0,0)$ is a P-position.
\end{theorem}

After any move of the first player that cuts an infinite column of $\PA$ at height $z$, the
second player cuts the poison column at height $z+1$ or $z+2$. The finite positions that arise
belong to explicit families (Theorem~\ref{thm:families}). Theorem~\ref{thm:A} is machine-checked
in Lean~4 (Section~\ref{sec:lean}).

\begin{theorem}\label{thm:B}
$\PB=(\w,\w,2\mid\w,\w,2\mid 2,2,0)$ is an N-position. A winning move cuts the column $(1,0)$ at
height $31$; no cut of that column at a lower height wins.
\end{theorem}

$\PB$ has no finite winning reply up to height $253$, and every cut of its column $(1,0)$ at a
height below $31$ has a winning reply. So a search of bounded depth sees no difference between
$\PA$ and $\PB$. Section~\ref{sec:chains} explains what happens at height $31$.

\begin{theorem}\label{thm:class}
Consider the positions whose finite columns have height at most $4$ and whose set of
infinite columns is one of
\begin{itemize}
\item the row of three $\{(0,0),(1,0),(2,0)\}$,
\item the tripod $\{(0,0),(0,1),(1,0)\}$,
\item the row of three with one more cell, $\{(0,0),(1,0),(2,0),(0,1)\}$,
\item the $2\times2$ block.
\end{itemize}
There are $490+1001+315+371=2177$ of them. The only P-positions among them are $\PA$
and its transpose. Each of the others has a winning move to a finite P-position or to a
P-position with two infinite columns.
\end{theorem}

Up to transposition these are all the sets of three or four infinite columns over the
$3\times3$ base.

\begin{theorem}\label{thm:cross}
Every position with the five columns $(0,0),(0,1),(0,2),(1,0),(2,0)$ infinite and the four
others of height at most $8$ has a winning move that cuts the poison column, at height at
most $70$. There are $825$ such positions. The same holds, with height at most $23$, for the
$1001$ tripod positions of Theorem~\ref{thm:class}.
\end{theorem}

The towers of $3\times3\times\omega$ obtained by cutting the column $(1,1)$ are of the first
kind.

Theorems~\ref{thm:A}--\ref{thm:cross} are computer-assisted; Theorem~\ref{thm:A} is in
addition fully formalised. Section~\ref{sec:exact} says precisely what is computed and how it
was checked. Section~\ref{sec:law} reports an empirical
law for the heights at which chains such as the one in Theorem~\ref{thm:B} get their P member,
and Section~\ref{sec:open} lists what remains open. The programs and their outputs are in the
repository~\cite{repo}.

\section{Positions, and what is computed exactly}\label{sec:exact}

\subsection{Positions}
A \emph{position} is a function $a\colon\{0,1,2\}^2\to\mathbb N\cup\{\w\}$ that is weakly
decreasing in both coordinates, with $a_{00}\ge1$. It stands for the set of cells $(x,y,z)$ with
$z<a_{xy}$. A \emph{bite} at $(x,y,t)$, where $t<a_{xy}$ and $(x,y,t)\ne(0,0,0)$, replaces
$a_{x'y'}$ by $\min(a_{x'y'},t)$ for all $x'\ge x$, $y'\ge y$. We also say that it \emph{cuts}
the column $(x,y)$ at height $t$, and write $\cut{(x,y)}{t}$. The player who cannot bite
loses, so the position with $a_{00}=1$ and all other heights $0$ is a P-position. Every play is
finite, because the positions are well-founded under inclusion. The \emph{transpose} of a
position exchanges the two base coordinates; it has the same outcome.

The column $(0,0)$ is the \emph{poison column}. An infinite column can only stand on cells
whose predecessors are infinite, so the infinite columns form a down-set of the base that
contains $(0,0)$.

\begin{lemma}[chains]\label{lem:chain}
Let $X$ be a position and $(x,y)$ a column. Among the positions obtained from $X$ by cutting
$(x,y)$ at the heights $t=0,1,2,\dots$ (with $t\ge1$ for the poison column), together with $X$
itself, at most one is a P-position.
\end{lemma}
\begin{proof}
The position cut at $t$ is an option of the position cut at $t'>t$ and of $X$.
\end{proof}

\subsection{Finite positions}
For a finite position write $r$ for the heights of the eight columns other than the poison
column. By Lemma~\ref{lem:chain} at most one height $a$ of the poison column makes $[a;r]$ a
P-position; call it $\pi(r)$. It is determined by the values of $\pi$ on smaller arguments:
no $[a;r]$ is P if $\pi(\min(r,a'))=a'$ for some $1\le a'<\max r$, and otherwise $\pi(r)$ is
the least $a\ge\max(r,1)$ that is not $\pi(r')$ for an $r'$ obtained from $r$ by one bite.
The program \texttt{base3\_table} tabulates $\pi$ for all $r$ within given bounds on the
columns, up to height $253$. It was compared with a brute-force search on $33$ settings of
small height, and it reproduces the known P-heights of small boxes and all values obtained
independently by the solver below. By \emph{a finite P cut} of a position we mean a cut of
its poison column that leaves a P-position; it is a winning move to a finite position.

\subsection{One infinite column}
If only the poison column is infinite, the position is an N-position: with $r$ the other
heights, either $\pi(r)$ exists and the cut of the poison column at $\pi(r)$ wins, or some
lower cut does \cite{HS,Gale}.

\subsection{Two infinite columns}\label{sec:two}
Let the poison column and one neighbouring column, say $(0,1)$, be infinite, with finite
heights $r$ elsewhere. Cutting $(0,1)$ leaves one infinite column and never wins. So the
position is P if and only if no bite in a finite column leaves a P-position and no cut of the
poison column does. The first condition refers to finitely many positions of the same kind
with smaller $r$. The second concerns the sequence $p\mapsto\pi(p,r)-p$, where $p$ is the
height of $(0,1)$: the cut at $p$ is P exactly when this is $0$. For $p$ beyond the entries of
$r$ this sequence, taken jointly for all $r'\le r$, is produced by a deterministic machine
with a bounded state, so it is eventually periodic (compare the periodicity theorem of
Byrnes \cite{By}). The solver \texttt{rodf} simulates the machine until a state repeats. A
repeated state proves that the sequence never vanishes, or exhibits the height where it does.
We call a position with two infinite columns \emph{certified} P (or N) when the solver has
closed these certificates for it and, recursively, for every position its verdict depends on.
Certified sequences in this paper have up to $5.5\cdot10^6$ terms before the state repeats.

The solver has a mode without certificates in which the search is cut off at a height. Such
verdicts are not used.

\subsection{What the theorems rest on}
\begin{itemize}
\item Theorem~\ref{thm:A}, as proved in Section~\ref{sec:A}: an SMT computation (Z3 \cite{Z3},
repeated with cvc5) for two statements of Presburger arithmetic; a hand argument
(Lemma~\ref{lem:W}); three certified positions with two infinite columns; and the fact that
the rod $2\times\omega$ is a P-position \cite{HS}. Independently of all of these, the theorem
is proved in Lean~4 (Section~\ref{sec:lean}); that proof rests on the Lean kernel only.
\item Theorems~\ref{thm:B} and~\ref{thm:class}: certified positions, and the table of finite
positions.
\item Theorem~\ref{thm:cross}: the table of finite positions.
\end{itemize}
A third program, the \emph{census} (Section~\ref{sec:class}), searches for winning replies. It
is used to find them and never to prove anything.

\section{The position \texorpdfstring{$\PA$}{A}}\label{sec:A}

The bottom layer of $\PA$ is the $3\times3$ square without the cells $(2,1)$ and $(2,2)$; the
second layer is the $2\times2$ block together with $(2,0)$; from height $2$ on only the block
remains.

\subsection{Finite positions below \texorpdfstring{$\PA$}{A}}
They live on seven columns. Write
\[
(a;\,b,c,d\mid e,f\mid k)
\]
for the heights of $(0,0)$; $(0,1)$, $(1,0)$, $(1,1)$; $(0,2)$, $(1,2)$; and $(2,0)$. So
$e,f\le1$ and $k\le2$. We omit a trailing $\mid0,0\mid0$. Four \emph{shapes} are used:
\[
\begin{aligned}
&\text{shape 1: } (m+1;\,m+1,m,m), &&\text{shape 2: } (m+1;\,m,m+1,m),\\
&\text{shape 3: } (m+2;\,m+2,m+2,m), &&\text{shape 4: } (m+2;\,m+1,m+1,m).
\end{aligned}
\]

\begin{theorem}\label{thm:families}
The following finite positions are P-positions.
\begin{center}
\begin{tabular}{lll}
\toprule
name & position & for\\
\midrule
$F_1$--$F_4$ & shapes 1--4 with $e=f=1$, $k=2$ & $m\ge2$ for $F_1$; $m\ge1$ otherwise\\
$K_1$--$K_4$ & shapes 1--4 with $e=1$, $f=0$, $k=1$ & $m\ge2$ for $K_1,K_2$; $m\ge1$ otherwise\\
$G(c,d)$ & $(c+2;\,d+2,c,d\mid1,1\mid0)$ & $c\ge d\ge1$\\
$T_1$ & $(m+2;\,1,m,1\mid1,0\mid1)$ & $m\ge2$\\
$T_2$ & $(m+2;\,m,1,1\mid1,0\mid1)$ & $m\ge2$\\
$T_3$ & $(m+3;\,m,1,1\mid1,1\mid1)$ & $m\ge1$\\
\bottomrule
\end{tabular}
\end{center}
\end{theorem}

In $G(c,d)$ the layer $x=1$ is the layer $x=0$ lowered by $2$ in both columns, and the second
player can keep it so: a cut of one layer at height $t$ is answered by a cut of the other at
$t\mp2$.

\begin{lemma}\label{lem:W}
$W(m)=(m+1;\,m,m+1,m\mid1,1\mid0)$ is an N-position for every $m\ge1$.
\end{lemma}

$W(m)$ is shape 2 with $e=f=1$ and $k=0$: the three-row position $[m+1,m,1]$ taken twice, side
by side. Its winning replies show no pattern (Remark~\ref{rem:W}); the proof does not find
them.

\subsection{The closed set}
Let $S$ be the union of the positions of Theorem~\ref{thm:families} and of
\begin{itemize}
\item \emph{pure}: the positions with $e=f=k=0$ and either $b+c=a-1$, $d=\min(b,c)$, or
$(a;\,a,a,a-1)$. These are positions on the $2\times2$ base alone;
\item \emph{flat}: positions on three columns in a line. The heights are $[2,2,1]$ or
$[3,1,1]$, on the line $y=0$ or on the line $x=0$, or $[c+2,c,2]$ with $c\ge2$ on the line
$y=0$ (the columns $(0,0)$, $(1,0)$, $(2,0)$);
\item \emph{small}: $106$ further positions of height at most $8$, listed in the script.
\end{itemize}
Let $S_0$ be $S$ without $F_1,\dots,F_4$.

\begin{claim}\label{cl:S0}
For every $X\in S_0$ and every option $Y$ of $X$, some option of $Y$ is in $S_0$.
\end{claim}
\begin{claim}\label{cl:S}
For every $X\in S$ and every option $Y$ of $X$: either $Y=W(m)$ for some $m\ge1$, or some
option of $Y$ is in $S$.
\end{claim}

A bite replaces some heights by their minimum with $t$, and each family is given by linear
equations and inequalities. So both claims are sentences of Presburger arithmetic. The script
\texttt{candidate\_A\_closure\_z3.py} proves them: for each of the $14$ families and each of
the seven columns, Z3 refutes the statement ``some $X$ in the family and some height $t$ give
an option that is not a $W(m)$ and has no option in the set''. The reply is existentially
quantified inside the refuted formula. The small positions are checked by enumeration. The
whole run takes about one minute.

Three checks were made on the script. The same predicates, evaluated on integers, confirm both
claims for all positions of height at most $14$. Without the exception for $W(m)$ the solver
returns the expected counterexample ($F_2$ with the column $(2,0)$ removed), and without the
family $G$ another one ($F_3$ with that column removed). All $4404$ positions of $S$ of height
at most $60$ are P-positions in the table of finite positions, and no $W(m)$ with $m\le117$ is.

\begin{proof}[Proof of Theorem~\ref{thm:families} for $S_0$]
By induction on the number of cells. Let $X\in S_0$. By Claim~\ref{cl:S0} every option of $X$
has an option in $S_0$, which is smaller and hence P by induction. So every option of $X$ is N.
\end{proof}

\subsection{Proof of Lemma~\ref{lem:W}}
For $m=1$ the position $W(1)$ has an option in $S_0$. Let $m\ge2$ and let $u=(1,1,m-1)$ be the
top cell of the column $(1,1)$. It is a maximal cell of $W(m)$: the only column above $(1,1)$
in the base that is not empty is $(1,2)$, of height $1\le m-1$. So
$W'=W(m)\setminus\{u\}=(m+1;\,m,m+1,m-1\mid1,1\mid0)$ is an option of $W(m)$.

Suppose $W(m)$ is a P-position. Then $W'$ is N, so some bite at a cell $q$ of $W'$ leaves a
P-position. If $q\le u$, the same bite on $W(m)$ removes $u$ as well and leaves the same
position, which is then a P option of $W(m)$: a contradiction. If $q\not\le u$, then $q$ lies
in one of the columns $(0,2)$, $(1,2)$ or has height at least $m$; only the columns $(0,0)$ and
$(1,0)$ reach height $m+1$. So there are four such cells, and none leaves a P-position:
\begin{center}
\begin{tabular}{lll}
\toprule
$q$ & position left & an option of it in $S_0$\\
\midrule
$(0,0,m)$ & $(m;\,m,m,m-1\mid1,1\mid0)$ & take $(0,2,0)$: pure $(m;\,m,m,m-1)$\\
$(1,0,m)$ & $(m+1;\,m,m,m-1\mid1,1\mid0)$ & take $(0,1,0)$: pure $(m+1;\,0,m,0)$\\
$(0,2,0)$ & $(m+1;\,m,m+1,m-1)$ & take $(1,0,0)$: pure $(m+1;\,m,0,0)$\\
$(1,2,0)$ & $(m+1;\,m,m+1,m-1\mid1,0\mid0)$ & cut $(0,0)$ at $1$: a small position\\
\bottomrule
\end{tabular}
\end{center}
The options in the last column are P by the part of Theorem~\ref{thm:families} already proved,
so the positions in the middle column are N. This contradiction shows that $W(m)$ is N. \qed

\medskip
The argument is strategy stealing on a single cell: $W(m)$ without its top $(1,1)$ cell has
only four options that $W(m)$ does not have, and each of them loses for a simple reason.

\begin{proof}[Proof of Theorem~\ref{thm:families}]
By induction on the number of cells. Let $X\in S$. By Claim~\ref{cl:S} every option $Y$ of $X$
is a $W(m)$, which is N by Lemma~\ref{lem:W}, or has an option in $S$, which is P by induction.
\end{proof}

\subsection{Proof of Theorem~\ref{thm:A}}
The options of $\PA$ are the bites in its four infinite columns and in its three finite ones.
Each has a P option:
\begin{center}
\begin{tabular}{lll}
\toprule
option of $\PA$ & reply & the reply leaves\\
\midrule
cut $(0,0)$ at $h\ge3$ & cut $(1,0)$ at $h-1$ & $F_1(h-1)$\\
cut $(0,0)$ at $2$ & take $(1,0,0)$ & flat $[2,2,1]$\\
cut $(0,0)$ at $1$ & take $(1,1,0)$ & $(1;\,1,1,0\mid1,0\mid1)$, small\\
cut $(1,0)$ at $z\ge2$ & cut $(0,0)$ at $z+1$ & $F_1(z)$\\
cut $(1,0)$ at $1$ & take $(1,2,0)$ & $(\w,\w,1\mid1,1,0\mid1,0,0)$, certified\\
cut $(1,0)$ at $0$ & cut $(0,0)$ at $2$ & flat $[2,2,1]$\\
cut $(0,1)$ at $z\ge1$ & cut $(0,0)$ at $z+1$ & $F_2(z)$\\
cut $(0,1)$ at $0$ & take $(2,0,0)$ & the rod $2\times\omega$ on $(0,0),(1,0)$\\
cut $(1,1)$ at $z\ge1$ & cut $(0,0)$ at $z+2$ & $F_3(z)$\\
cut $(1,1)$ at $0$ & cut $(0,0)$ at $1$ & $(1;\,1,1,0\mid1,0\mid1)$, small\\
take $(0,2,0)$ & take $(1,0,0)$ & the rod $2\times\omega$ on $(0,0),(0,1)$\\
take $(1,2,0)$ & cut $(0,1)$ at $2$ & $(\w,2,1\mid\w,2,0\mid2,0,0)$, certified\\
take $(2,0,1)$ & cut $(0,1)$ at $3$ & $(\w,3,1\mid\w,3,1\mid1,0,0)$, certified\\
take $(2,0,0)$ & take $(0,1,0)$ & the rod $2\times\omega$ on $(0,0),(1,0)$\\
\bottomrule
\end{tabular}
\end{center}
The positions $F_1(z)$ for $z\ge2$, $F_2(z)$ and $F_3(z)$ for $z\ge1$, and the flat and small
ones are P by Theorem~\ref{thm:families}. The rod $2\times\omega$ is P. The three remaining
positions have two infinite columns and are certified P. So every option of $\PA$ is N. \qed

\subsection{A machine-checked proof}\label{sec:lean}
Theorem~\ref{thm:A} is proved in Lean~4, with the core library only, as the statement
\texttt{Chomp.posA\_isP}: the set of cells
\[
\{(x,y,z):x\le1,\ y\le1\}\ \cup\ \{(x,2,0):x\le1\}\ \cup\ \{(2,0,z):z\le1\}
\]
of $\mathbb N^3$ is a P-position. Here N- and P-positions are defined inductively from the
rule that a move removes a cell other than $(0,0,0)$ together with all cells above it. The
proof depends on the axioms \texttt{propext}, \texttt{Classical.choice} and
\texttt{Quot.sound} and on nothing else; in particular it uses no SMT solver, no table and no
certificate from our solver. It follows the proof above, with three differences.

\emph{Explicit replies.} Claims~\ref{cl:S0} and~\ref{cl:S} assert that replies exist. In Lean
each reply is given: for every family and every column there is a decision tree over
conditions such as $t<3$, $t+2=c$ or ``$t$ is even'', and each leaf names the column and the
height of the reply and the family it lands in. Lean checks every leaf by linear arithmetic.
The trees are found by a script that searches small cases and asks Z3 for cases not yet
covered; Z3 only proposes. The small positions are checked by evaluation.

\emph{The three certified positions.} They are proved to be P-positions in Lean. The
positions that answer their finite cuts are not all in $S$, so $S$ is enlarged by the
families
\begin{center}
\begin{tabular}{lll}
\toprule
position & for even $c$ & for odd $c$\\
\midrule
$(a;\,1,c,1\mid1,0\mid2)$ & $a=c+3$, $c\ge2$ & $a=c+1$, $c\ge3$\\
$(a;\,1,c,1\mid1,1\mid1)$ & $a=c+3$, $c\ge2$ & $a=c+3$, $c\ge1$\\
$(a;\,2,c,1\mid1,0\mid1)$ & $a=c+1$, $c\ge2$ & $a=c+3$, $c\ge1$\\
$(a;\,2,c,1\mid1,1\mid1)$ & $a=c+2$, $c\ge2$ & $a=c+4$, $c\ge1$\\
$(a;\,2,c,2\mid1,0\mid2)$ & $a=c+1$, $c\ge2$ & $a=c+3$, $c\ge3$\\
$(a;\,3,c,3\mid1,1\mid1)$ & $a=c+2$, $c\ge6$ & $a=c+4$, $c\ge7$\\
\bottomrule
\end{tabular}
\end{center}
and by the four positions $(9;\,3,4,3\mid1,1\mid1)$, $(8;\,2,4,2\mid1,1\mid1)$,
$(5;\,3,4,1\mid1,1\mid1)$ and $(9;\,3,3,3\mid1,1\mid0)$. The second line of the table is the
transpose of $T_3$. The enlarged set, $26$ families and $114$ small positions, is again closed
up to the positions $W(m)$: $182$ cases, one for each family and column. The proof involves
$46$ positions with two infinite columns; seven of them are P-positions
(the three, the transpose of the first, the position $(\w,2,1\mid\w,1,1\mid1,0,0)$, and the rods
$2\times\omega$ in both directions), and each of the others is N by one explicit move.

\emph{One infinite column.} The Fundamental Theorem is proved for the positions at hand: if
only the poison column is infinite, then for each finite part at most one height of the
poison column gives a P-position, so the heights that are answered by a bite elsewhere are
bounded, and above that bound a cut of the poison column must win.

The hand-written part is about $600$ lines and the generated part about $3800$; the whole
library compiles in under three minutes.

\begin{remark}\label{rem:W}
The winning replies to $W(m)$ are of two kinds, a cut of $(0,1)$ at $t_1$ or a cut of $(1,1)$
at $t_2$. Some values (a dash means that no reply of that kind wins):
\begin{center}
\small\setlength{\tabcolsep}{4pt}
\begin{tabular}{l*{17}{c}}
\toprule
$m$ & 3 & 4 & 5 & 6 & 7 & 8 & 9 & 10 & 11 & 12 & 17 & 19 & 22 & 30 & 44 & 60 & 80\\
\midrule
$t_1$ & 1 & 2 & 3 & 4 & 5 & -- & 7 & 8 & 9 & 10 & -- & 16 & -- & -- & 35 & -- & --\\
$t_2$ & -- & 1 & 1 & 1 & 4 & 4 & 4 & 4 & 8 & 8 & 8 & -- & 16 & 24 & -- & 41 & 67\\
\bottomrule
\end{tabular}
\end{center}
The first kind is a matching with
gaps, like the partner function of three-row Chomp \cite{BHMS,Z}. A direct proof of
Theorem~\ref{thm:A} along the lines of the table would need a statement about this matching
for all $m$. Lemma~\ref{lem:W} avoids it.
\end{remark}

\begin{corollary}
The positions that have $\PA$ as an option are N-positions. They are
$(\w,\w,1\mid\w,\w,1\mid k,0,0)$ with $3\le k\le\w$; $(\w,\w,j\mid\w,\w,i\mid2,0,0)$ with
$2\le j\le\w$ and $1\le i\le j$; and $(\w,\w,1\mid\w,\w,1\mid2,e,f)$ with $1\le e\le2$,
$f\le1$.
\end{corollary}

\section{Chains of cuts, and the position \texorpdfstring{$\PB$}{B}}\label{sec:chains}

Let $X$ be a position whose infinite columns are exactly the $2\times2$ block. For $v\le u$ let
\[
S(u,v)=X\text{ after cutting }(1,0)\text{ at }u\text{ and }(1,1)\text{ at }v .
\]
The positions $S(u,u)$, $u=0,1,2,\dots$, are the options of $X$ by a cut of $(1,0)$; we call
them the \emph{chain} of $(1,0)$. By Lemma~\ref{lem:chain} at most one of them is P, and $X$ is
N if one is. Every $S(u,v)$ has exactly two infinite columns, $(0,0)$ and $(0,1)$, so its
outcome can be certified.

Call the cell $(u,v)$ \emph{eligible} if $S(u,v)$ has no finite P cut and no P option by a bite
in one of the columns $(0,2),(1,2),(2,0),(2,1),(2,2)$.

\begin{proposition}[the matching]\label{prop:matching}
Suppose that no $S(t,t)$ with $t<v$ is a P-position. Then $S(u,v)$ is a P-position if and only
if $(u,v)$ is eligible, no $S(t,v)$ with $v\le t<u$ is a P-position, and no $S(u,t)$ with
$t<v$ is a P-position.
\end{proposition}
\begin{proof}
The options of $S(u,v)$ are: the cuts of the poison column; the cuts of $(0,1)$, which leave
one infinite column and are N; the bites in the five columns named above; the cuts of $(1,0)$
at $t<u$, which leave $S(t,v)$ if $t\ge v$ and $S(t,t)$ if $t<v$; and the cuts of $(1,1)$ at
$t<v$, which leave $S(u,t)$. The position is P if and only if none of these is P.
\end{proof}

So the P-positions among the $S(u,v)$ are found row by row: \emph{row $u$ takes the smallest
$v<u$ such that $(u,v)$ is eligible and no earlier row has taken $v$}. Call this $v$ the
\emph{partner} of $u$. Each row takes at most one column and each column is taken at most once.

\begin{corollary}\label{cor:chain}
\begin{enumerate}[label=\textup{(\roman*)}]
\item The first P-position in the chain of $(1,0)$ is $S(u,u)$ for the least $u$ such that
$(u,u)$ is eligible and row $u$ has no partner. If there is no such $u$, the chain has no
P-position.
\item Up to that $u$, the number of columns below $u$ that no row has taken increases by
exactly one at each row without a partner, and is otherwise constant.
\item If $(u,u)$ is not eligible for any $u\ge u_0$ and no $S(u,u)$ with $u<u_0$ is P, the
chain has no P-position.
\end{enumerate}
\end{corollary}

For $\PA$, case (iii) applies to both chains with $u_0=2$: the diagonal position $S(u,u)$ has the
finite P cut at height $u+1$, which leaves $F_1(u)$ (or $F_2(u)$ for the chain of $(0,1)$).
We call such a chain \emph{regular}.

A column can serve only finitely many rows if the position $X$ cut at $\cut{(1,1)}{v}$, which
has the tripod of infinite columns, has a finite P cut, say at height $k$: then $S(u,v)$ has
the same finite P cut for every $u\ge k$. In the cases computed this $k$ exists, and $k-v$
grows roughly in proportion to $v$. Below that height the eligibility of the cells of a
column shows no pattern. Two further constraints are easy to prove: for a fixed $h$, the
cells $(u,v)$ with $u<h$ for which $S(u,v)$ has its finite P cut at $h$ contain at most one
cell of each row and of each column; and the cells with a P option by a bite in one of the
five other columns contain at most $B$ cells of each row, where $B$ is the number of such
bites of $X$.

\subsection{Proof of Theorem~\ref{thm:B}}
For $X=\PB$ we have $S(u,v)=(\w,\w,2\mid u,v,2\mid 2,2,0)$ for $u\ge v\ge2$. The solver
certifies that $S(31,31)$ is a P-position; the verdict rests on $1559$ certified positions with
two infinite columns. This proves that $\PB$ is N. That no $S(u,u)$ with $u<31$ is P follows
from Lemma~\ref{lem:chain}. \qed

\medskip
The partners of the rows $u\ge5$, according to the census of Section~\ref{sec:class} (the
individual entries of this table are not certified):
\begin{center}
\small
\begin{tabular}{l*{14}{c}}
\toprule
$u$ & 5 & 6 & 7 & 8 & 9 & 10 & 11 & 12 & 13 & 14 & 15 & 16 & 17 & 18\\
partner & 4 & 5 & 6 & 7 & 8 & -- & 10 & 11 & 9 & 12 & 13 & 14 & 16 & 15\\
\midrule
$u$ & 19 & 20 & 21 & 22 & 23 & 24 & 25 & 26 & 27 & 28 & 29 & 30 & 31 &\\
partner & 17 & 18 & 19 & 20 & 21 & 22 & 24 & 25 & 23 & 26 & 28 & 27 & -- &\\
\bottomrule
\end{tabular}
\end{center}
Row $10$ has no partner, but its diagonal cell is not eligible: $S(10,10)$ has a finite P
cut at height $17$. Row $31$ has no partner either: the cells $(31,v)$ with $v=28,29,30$ have
finite P cuts, and the eligible ones, $v=22,23,24,25,27$, are taken. Its diagonal cell is
eligible. Every other row from $5$ to $30$ had one or two columns to choose from.

\section{Classification of small positions}\label{sec:class}

\subsection{The census}
To find winning replies we use a search over a region of positions: fix a height $H$, a bound
$Z<H$, and bounds on the finite columns; the columns designated as free take the values
$0,\dots,Z$ and $\w$. Going up from the smallest positions, a position with at least two infinite
columns is \emph{labelled P} if no cut of its poison column at a height up to $H$ is P and no
option is labelled P, where an infinite column is cut only at heights up to $Z$. The label is
the true outcome up to these two truncations. It can be wrong: we found positions with two
infinite columns and finite columns of height at most $8$ and $11$ whose only finite P cuts
are at heights $381$ and $978$, and which a census with $H=100$ and $H=240$ labels P. So the
census only proposes replies. A reply that leaves two infinite columns is then certified by
the solver of Section~\ref{sec:two}, and a reply to a finite position is exact.

\subsection{Proof of Theorem~\ref{thm:class}}
For each of the four sets of infinite columns the table of finite positions gives the
positions with a finite P cut. For each of the remaining positions other than $\PA$ and its
transpose the census proposes a cut of an infinite column that leaves two infinite columns,
and the solver certifies the position it leaves as P, or else shows that a lower cut of the
same column leaves a certified P-position. In two cases the reply was found by the solver's
own search. The counts are:
\begin{center}
\begin{tabular}{lrrrr}
\toprule
infinite columns & positions & P & N by a finite P cut & N by a certified reply\\
\midrule
row of three & 490 & 0 & 243 & 247\\
tripod & 1001 & 0 & 1001 & 0\\
$\{(0,0),(1,0),(2,0),(0,1)\}$ & 315 & 0 & 311 & 4\\
$2\times2$ block & 371 & 2 & 156 & 213\\
\bottomrule
\end{tabular}
\end{center}
The census had the right verdict for all $464$ positions that needed a reply and the wrong
height for $7$ of them. The replies are listed in the repository. $\PA$ is P by
Theorem~\ref{thm:A}. \qed

\medskip
For the row of three the result agrees with earlier computations in which all such positions
with finite columns up to height $8$ were decided.

\begin{remark}\label{rem:five}
The same method with finite columns of height at most $5$: of the $812$ positions with the
$2\times2$ block infinite, $362$ have a finite P cut, $444$ are N by a certified reply, $\PA$
and its transpose are P, and four are undecided, namely
$(\w,\w,5\mid\w,\w,5\mid5,2,2)$, $(\w,\w,5\mid\w,\w,5\mid4,3,3)$ and their transposes. For
these the census proposes replies at heights $68$ and $42$, which the solver does not reach;
every cut of $(1,0)$ or $(0,1)$ at a height up to $24$ is certified N. Section~\ref{sec:law}
explains why a few such positions are to be expected.
\end{remark}

\subsection{Proof of Theorem~\ref{thm:cross}}
The table of finite positions with the corner columns bounded by $8$ and the others by $64$
gives a finite P cut for $823$ of the $825$ positions, at height at most $57$ (at most $10$ for
half of them). The two others are $(\w,\w,\w\mid\w,7,7\mid\w,5,2)$ and its transpose; a table
for that position up to height $96$ gives the finite P cut at height $70$. For the tripod
positions the table up to height $200$ gives a finite P cut at height at most $23$. \qed

\medskip
Thirteen of the $105$ cross positions with corner columns up to height $4$ have their finite P
cut below their tallest corner column, for example $(\w,\w,\w\mid\w,3,2\mid\w,2,1)$ at height
$2$. The cuts get higher as the corner grows ($27$ for corner columns up to $4$, then $57$ and
$70$), so a table of bounded height shows apparent exceptions.

\subsection{Two rows of three}
With the $3\times2$ block $\{0,1,2\}\times\{0,1\}$ infinite and a \emph{roof}
$(f_0,f_1,f_2)$ on the row $y=2$, the cut of $(1,0)$ again leaves two infinite columns. For the
$35$ roofs with $f_0\le4$:
\begin{itemize}
\item \emph{N by a finite P cut} ($12$ roofs): $(2,2,0)$, $(2,1,1)$, $(2,2,1)$, $(3,2,1)$,
$(3,2,2)$, $(3,3,1)$, $(3,3,2)$, $(4,2,1)$, $(4,3,1)$, $(4,4,1)$, $(4,4,3)$, $(4,4,4)$.
\item \emph{N by a cut of $(1,0)$ that leaves a certified P-position} ($11$ roofs; the height
of the cut after the arrow): $(0,0,0)\to0$, $(1,0,0)\to2$, $(1,1,0)\to1$, $(1,1,1)\to4$,
$(2,0,0)\to18$, $(3,1,0)\to7$, $(3,1,1)\to1$, $(3,3,0)\to3$, $(3,3,3)\to1$, $(4,1,1)\to2$,
$(4,4,0)\to9$.
\item \emph{Undecided} ($12$ roofs; after the colon the height up to which every cut of
$(1,0)$ is certified N): $(2,1,0)$: $24$; $(2,2,2)$: $24$; $(3,0,0)$: $42$; $(3,2,0)$: $20$;
$(4,0,0)$: $20$; $(4,1,0)$: $20$; $(4,2,0)$: $20$; $(4,3,0)$: $19$; $(4,2,2)$: $17$;
$(4,3,2)$: $16$; $(4,3,3)$: $16$; $(4,4,2)$: $16$.
\end{itemize}
The roofs $(g,g,g)$ are the towers of $3\times3\times\omega$ obtained by cutting $(0,2)$ at
height $g$; so those with $g=1,3,4$ are N and the one with $g=2$ is undecided. The roof
$(2,0,0)$ has $\PA^{\mathsf T}$ as an option, so it is N also by Theorem~\ref{thm:A}.

\section{How chains end: an empirical law}\label{sec:law}

This section reports measurements. Nothing in it is proved about Chomp; the one
proposition concerns a model.

Take the $2\times2$ block with finite columns of height at most $4$, and for each of the $215$
positions without a finite P cut both chains (of $(1,0)$ and of $(0,1)$). This set of
positions is symmetric under transposition, which exchanges the two chains, so every chain
occurs twice; the counts below are not halved, and their statistical weight is half of what
they suggest. Follow each chain
from height $5$ (or the height of its tallest finite column) to its first P-position, as given
by the census with $H=253$, $Z=120$. Six chains are regular in the sense of
Corollary~\ref{cor:chain}(iii): the four of $\PA$ and its transpose, and one chain each of
$(\w,\w,1\mid\w,\w,1\mid0,0,0)$ and its transpose, where the diagonal position has its finite P
cut at height $u+2$ (it leaves $G(u,u)$). The other chain of these two positions has a
P-position at height $0$. For the $216$ other chains that reach height $5$:
\begin{center}
\begin{tabular}{lrrrrr}
\toprule
heights & rows & diagonal eligible & no partner & chain ends & ends${}\times{}$height\\
\midrule
5--9 & 794 & 33\% & 43\% & 122 & 0.98\\
10--19 & 644 & 50\% & 14\% & 50 & 1.03\\
20--39 & 584 & 47\% & 7.2\% & 24 & 1.10\\
40--69 & 426 & 42\% & 2.3\% & 8 & 0.97\\
70--99 & 262 & 44\% & 3.8\% & 4 & 1.26\\
100--120 & 168 & 44\% & 0\% & 0 & 0\\
\bottomrule
\end{tabular}
\end{center}
The last column is the number of chain ends divided by the sum of $1/u$ over the rows.

\begin{observation}\label{obs:law}
A chain that has no P-position below height $u$ has one at height $u$ with probability about
$G/u$, with $G=1.00$ ($208$ ends, that is $104$ up to transposition). Hence a chain without a
P-position below $s$ has none below $h$ with probability about $s/h$.
\end{observation}

The numbers of chains still without a P-position at the heights $10$, $20$, $40$, $70$,
$100$, $120$ are $82$, $42$, $20$, $12$, $8$, $8$; the law gives $86.1$, $43.0$, $21.4$,
$12.2$, $8.6$, $7.1$. The constant was first fitted
on a census with $H=160$, $Z=70$, which gave $G=1.01$ and predicted $8.4$ and $7.0$ for the
heights $100$ and $120$ that it did not cover.

The same count for finite columns of height at most $5$ ($812$ positions with the block
infinite, of which $450$ have no finite P cut; census with $H=160$, $Z=70$) gives $G=0.98$
with $534$ ends ($267$ up to transposition).

\subsection*{The mechanism}
Call a column $v<u$ \emph{live} in row $u$ if no $S(t,v)$ with $t\le u$ has a finite P cut
at a height at most $t$, and call the live columns that no earlier row has taken the
\emph{pool} of row $u$. By Corollary~\ref{cor:chain}(ii) the pool gains a column at each row
without a partner (and loses the columns that die unused). In the data, a row
at height at least $10$ whose pool has $k$ columns has no partner in $38\%$, $18\%$, $6.9\%$,
$4.6\%$, $0.7\%$, $0\%$ of the cases for $k=1,\dots,6$: close to $\delta^k$ with
$\delta=0.425$. On the rows without a partner the diagonal cell is eligible in $86$ cases of
$154$, so with frequency $\varepsilon=0.56$. The mean size of the pool grows from $2.8$ at
heights $10$--$19$ to $5.0$ at heights $100$--$120$.

Consider the model in which a row with pool $k$ has no partner with probability $\delta^k$,
independently of the past, and then ends the chain with probability $\varepsilon$ and
otherwise enlarges the pool by one.

\begin{proposition}\label{prop:model}
In this model let $T$ be the height at which the chain ends, and let
$\gamma=\log\frac1{1-\varepsilon}\big/\log\frac1\delta$. Then there are constants $0<c<C$,
depending on $\delta$, $\varepsilon$ and the initial pool, such that for all $h\ge1$
\[
c\,h^{-\gamma}\le\Pr(T>h)\le C\,h^{-\gamma}.
\]
\end{proposition}
\begin{proof}
The number $L_k$ of rows spent with pool $k$ is geometric, $\Pr(L_k>t)=(1-\delta^k)^t$. From
the initial pool $k_0$, the $m$-th row without a partner comes at
$\tau_m=L_{k_0}+\dots+L_{k_0+m-1}$, and the chain survives each such row with probability
$1-\varepsilon$. So $\Pr(T>h)=\mathbb E\,(1-\varepsilon)^{N(h)}$, where $N(h)$ is the number
of $m$ with $\tau_m\le h$. Let $\delta^{-(n-1)}<h\le\delta^{-n}$. Since
$\Pr(L_{n+1}>h)\ge(1-\delta/h)^h$ is bounded below and $L_{n+1}>h$ implies
$N(h)\le n-k_0+1$, we get $\Pr(T>h)\ge c'(1-\varepsilon)^{n}$. In the other direction,
$N(h)<n-k_0-j$ means $\tau_m>h$ for $m=n-k_0-j$, a sum of geometric variables whose means
are at most $\delta^{j}h,\ \delta^{j+1}h,\dots$; giving the term $r$ places below the last
the share $(1-\theta)\theta^rh$ of $h$, with $\delta<\theta<1$, shows that this has
probability at most $C''\exp(-c''\delta^{-j})$. Summing over $j$ gives
$\Pr(T>h)\le C(1-\varepsilon)^{n}$, and $(1-\varepsilon)^n$ is $h^{-\gamma}$ up to a constant
factor.
\end{proof}

With the measured $\delta$ and $\varepsilon$ the exponent is $\gamma=0.96$; for finite columns
up to height $5$ one finds $\delta=0.485$, $\varepsilon=0.527$ and $\gamma=1.03$. So the model
reproduces Observation~\ref{obs:law} from two frequencies measured separately from it. One has
$\gamma=1$ exactly when $\delta=1-\varepsilon$, that is, when the diagonal cell is eligible as
often as a cell of the pool: then a row has $k+1$ candidate cells of the same kind, takes the
first eligible one, and the chain ends when that is the last. In the data the diagonal cell
is eligible $46\%$ of the time and a pool cell $53\%$. The cells of a row are not independent
(their number of eligible cells is more concentrated than a binomial, as the first constraint
of Section~\ref{sec:chains} suggests), and there is no sign that chains differ durably from one
another: across the chains that reach height $40$, the frequencies of rows without a partner
at heights $10$--$19$ and $20$--$39$ have correlation $0.3$.

In the model every chain ends, with probability $1$. We have no proof of the corresponding
statement for Chomp (Section~\ref{sec:open}).

For the $3\times2$ block the same count on the $35$ roofs of Section~\ref{sec:class} gives
three chain ends between height $5$ and the end of the certified runs, where
Observation~\ref{obs:law} with $G=1$ would give $21$. Chains in that family end much more
rarely.

The law has two consequences for searches. A chain that has been followed to height $s$
without a P-position still has one below $2s$ with probability one half, so a bounded search
cannot separate a position like $\PA$ from one like $\PB$. And the only chains for which we have
a reason to expect no P-position at all are the regular ones. A scan of the $4917$ positions
with the $2\times2$ block infinite and finite columns up to height $8$, for those in which every
cut of an infinite column at a height from $14$ to $50$ has a finite winning reply, found
$\PA$, its transpose, and two symmetric positions that are N by a cut of $(0,1)$ at height $1$.

\section{Open questions}\label{sec:open}
\begin{enumerate}
\item Is $3\times3\times\omega$ an N-position? Its options obtained by cutting $(1,1)$ at a
height up to $14$ are N by finite P cuts as in Theorem~\ref{thm:cross}, and those obtained by
cutting $(0,1)$ at a height up to $11$ are N by replies that leave two infinite columns
\cite{repo}. Among the cuts of $(0,2)$, those at heights $1,3,4,6$ are N and the one at
height $2$ is undecided (Section~\ref{sec:class}). The cuts of $(1,2)$ at heights $1,2,5,11$
and of $(2,2)$ at heights $1,2,3,5,6,9,11$ are N by finite P cuts; for the other heights no
method is available, because every other option of such a position has at least three
infinite columns.
\item Is $(\w,\w,1\mid\w,\w,1\mid9,2,1)$ a P-position? (The same question arises for the four
positions of Remark~\ref{rem:five}.) It has no finite P cut up to height
$240$, the cuts of $(0,1)$ up to height $59$ and of $(1,0)$ up to height $26$ are certified N,
and the census finds no reply up to height $120$. Observation~\ref{obs:law} makes it an
expected survivor rather than a likely P-position.
\item Does every position with the cross or the tripod of infinite columns have a finite P
cut?
\item Prove that a chain whose diagonal cell is eligible infinitely often contains a
P-position. The constraints of Section~\ref{sec:chains} do not imply it: the pattern in which
every cell $(u,u-1)$ is eligible, the diagonal cell is eligible for even $u$, and every other
cell is not, satisfies them and never ends.
\item Is it true that a chain whose diagonal cell is eligible only finitely often has, from
some height on, the finite P cut of $S(u,u)$ at $u+c$ for a constant $c$? Together with the
previous statement this would reduce the outcome of a position with the $2\times2$ block
infinite to statements about finite positions in linear families, of the kind proved for
$\PA$.
\end{enumerate}

\appendix
\section{Reproducing the computations}
All programs are in the repository~\cite{repo}. The main ones:
\begin{itemize}
\item \texttt{scripts/candidate\_A\_closure\_z3.py}: Claims~\ref{cl:S0} and~\ref{cl:S} and the
four replies in the proof of Lemma~\ref{lem:W}.
\item \texttt{lean/Chomp/Block.lean}, \texttt{BlockFinite.lean} (written by hand) and
\texttt{BlockFamilies.lean}, \texttt{BlockA.lean} (generated by
\texttt{scripts/lean\_A\_gen.py}): the Lean proof of Theorem~\ref{thm:A}. Build with
\texttt{lake build}.
\item \texttt{tools/base3\_table.cpp}: the table of finite positions; with \texttt{-T} the
census.
\item \texttt{tools/rodf.cpp}: the solver for two infinite columns. The three certified
positions of Theorem~\ref{thm:A} and the one of Theorem~\ref{thm:B} are
\begin{quote}\ttfamily
rodf -p w,w,1,1,1,1 321\\
rodf -p w,2,1,w,2,2 321\\
rodf -p w,3,1,w,3,1,1 331\\
rodf -L 60000000 -p w,w,2,31,31,2,2,2 332
\end{quote}
\item \texttt{scripts/verify\_census\_options.py}: sends the replies proposed by the census
to the solver. Its outputs for Theorem~\ref{thm:class} are in \texttt{data/census/}.
\item \texttt{scripts/chain\_matching.py}, \texttt{scripts/chain\_survival.py}: the tables of
Sections~\ref{sec:chains} and~\ref{sec:law}.
\end{itemize}

\begin{thebibliography}{9}
\bibitem{BHMS} A.~E.~Brouwer, G.~Horv\'ath, I.~Moln\'ar-S\'aska, C.~Szab\'o, On three-rowed
  Chomp, \emph{Integers} \textbf{5} (2005), G07.
\bibitem{By} S.~Byrnes, Poset game periodicity, \emph{Integers} \textbf{3} (2003), G03.
\bibitem{Gale} D.~Gale, \emph{Tracking the Automatic Ant and Other Mathematical Explorations},
  Springer, 1998, ch.~11.
\bibitem{HS} S.~Huddleston, J.~Shurman, Transfinite Chomp, in \emph{More Games of No Chance},
  MSRI Publications \textbf{42}, Cambridge University Press, 2002, 183--212.
\bibitem{Now} R.~J.~Nowakowski, Unsolved problems in combinatorial games, in \emph{Games of
  No Chance 6}, MSRI Publications \textbf{71}, 2025.
\bibitem{repo} C.~Reams, Transfinite Chomp on $\omega\times\omega\times\omega$: programs, data and
  Lean proofs, \url{https://github.com/calr3/transfinite-chomp-wxwxw}, 2026.
\bibitem{Z3} L.~de Moura, N.~Bj{\o}rner, Z3: an efficient SMT solver, \emph{TACAS 2008},
  LNCS \textbf{4963}, 337--340.
\bibitem{Z} D.~Zeilberger, Three-rowed Chomp, \emph{Adv.\ Appl.\ Math.} \textbf{26} (2001),
  168--179.
\end{thebibliography}

\end{document}
